\documentclass[UTF8,zihao=-4]{ctexart}
\usepackage[a4paper,margin=2.1cm]{geometry}
\usepackage{hyperref}
\usepackage{array}
\usepackage{longtable}
\usepackage{verbatim}
\hypersetup{hidelinks}
\setlength{\parindent}{0pt}
\setlength{\parskip}{0.55em}
\emergencystretch=4em
\sloppy
\title{04-代码生成与优化复习讲义}
\author{}
\date{}
\begin{document}
\maketitle
\setcounter{tocdepth}{1}
\tableofcontents
\newpage
\section{04 代码生成与优化复习讲义}

对应资料：

\begin{itemize}
\item \texttt{Slides/Slides/Ch8-1-Target Code Generation.pdf}
\item \texttt{Slides/Slides/Ch8-2-Translation out of SSA.pdf}
\item \texttt{Slides/Slides/Ch8-3-Register Allocation.pdf}
\item \texttt{Slides/Slides/Ch10-Instruction Scheduling.pdf}
\item \texttt{Review/compile.pdf}
\item \texttt{Exam/2025.pdf} 指令调度题
\item \texttt{Exam/过时/2020-final-exam.pdf} 目标代码生成题型
\item \texttt{Exam/过时/2023-exam-A.pdf} 目标代码生成题型
\end{itemize}

\subsection{一、目标代码生成}

目标代码生成把 IR 转换为目标机器指令。核心任务：

\begin{itemize}
\item 指令选择。
\item 寄存器分配。
\item 指令调度。
\item 调用约定和栈帧处理。
\item 地址计算和内存访问。
\end{itemize}

\subsection{二、寻址模式}

常见模式：

\begin{quote}
\footnotesize
\begin{verbatim}
LD R1, a(R2)     : 取地址 (*R2)+a 处数据放入 R1
LD R1, 100(R2)   : 取地址 (*R2)+100 处数据放入 R1
LD R1, *R2       : 取地址 *R2 处数据放入 R1
LD R1, *100(R2)  : 先取 (*R2)+100 处数据作为新地址，再取该地址数据
LD R1, #100      : 将立即数 100 放入 R1
\end{verbatim}
\end{quote}

\subsection{三、块内优化}

常见优化：

\begin{itemize}
\item 局部公共子表达式消除。
\item 强度削弱。
\item 使用机器习语。
\item 去除冗余 LOAD/STORE。
\item 窥孔优化。
\end{itemize}

例：

\begin{quote}
\footnotesize
\begin{verbatim}
x * 2      -> x + x
x = x + 1  -> INC x
t1 = a+b; t2 = a+b -> t2 = t1
\end{verbatim}
\end{quote}

\subsection{四、从 SSA 转出}

phi 消除思路：

\begin{quote}
\footnotesize
\begin{verbatim}
a = phi(a1 from P1, a2 from P2)
\end{verbatim}
\end{quote}

改写：

\begin{quote}
\footnotesize
\begin{verbatim}
P1 末尾插入 a = a1
P2 末尾插入 a = a2
删除 phi
\end{verbatim}
\end{quote}

完整流程要处理并行复制：

\begin{enumerate}
\item 把 phi 的输入复制放到对应前驱块末尾。
\item 多个复制要按并行赋值语义处理。
\item 存在环形复制时引入临时变量。
\item 删除多余复制。
\end{enumerate}

\subsection{五、寄存器分配}

目标：

\begin{quote}
\footnotesize
\begin{verbatim}
无限虚拟寄存器 -> 有限物理寄存器
\end{verbatim}
\end{quote}

约束：

\begin{itemize}
\item 生成正确代码。
\item 使用不超过 k 个物理寄存器。
\item 减少 spill 的 load/store。
\item 控制栈空间。
\end{itemize}

\subsection{六、局部寄存器分配}

基本块内没有控制流分叉，可计算变量活跃区间。

\texttt{MAXLIVE}：某个程序点之后仍活跃的变量数。

若 \texttt{MAXLIVE > k}，要 spill。选择 spill 的经验规则：

\begin{quote}
\footnotesize
\begin{verbatim}
选择下一次使用距离最远的变量 spill。
\end{verbatim}
\end{quote}

\subsection{七、全局寄存器分配与图染色}

冲突图：

\begin{itemize}
\item 顶点：变量或虚拟寄存器。
\item 边：两个变量活跃区间相交，不能分配同一物理寄存器。
\end{itemize}

k-染色：

\begin{itemize}
\item 若冲突图可 k-染色，则可用 k 个寄存器完成分配。
\item 若不可染色，要 spill。
\end{itemize}

Chaitin 算法：

\begin{enumerate}
\item 找度数小于 k 的点，删除并压栈。
\item 重复直到图空。
\item 若找不到度数小于 k 的点，选一个点 spill。
\item 出栈时分配颜色。
\end{enumerate}

Chaitin-Briggs：

\begin{itemize}
\item 度数大于等于 k 的点先不直接 spill，也压栈。
\item 出栈染色失败时再 spill。
\end{itemize}

\subsection{八、指令调度}

目的：

\begin{itemize}
\item 把相互独立的指令放近，让硬件流水线和乱序执行利用并行性。
\item 减少停顿。
\item 提高硬件利用率。
\end{itemize}

依赖类型：

\begin{quote}
\footnotesize
\begin{verbatim}
True dependence / RAW  : 写后读。i 写 x，j 读 x，i 必须在 j 前。
Anti dependence / WAR  : 读后写。i 读 x，j 写 x，j 不能提前破坏 i 读到的值。
Output dependence / WAW: 写后写。i 写 x，j 写 x，二者顺序影响最终值。
\end{verbatim}
\end{quote}

寄存器重命名：

\begin{itemize}
\item 可消除 Anti dependence 和 Output dependence。
\item 不能消除 True dependence。
\end{itemize}

\subsection{九、列表调度}

步骤：

\begin{enumerate}
\item 建依赖图，节点是指令，边是依赖。
\item 计算每个节点的就绪条件。
\item 每个周期从 ready list 中选可发射指令。
\item 使用优先级打破并列，例如关键路径长度更大者优先。
\item 得到满足依赖的拓扑序。
\end{enumerate}

考场写法：

\begin{quote}
\footnotesize
\begin{verbatim}
依赖图 -> ready list -> 每周期选择 -> 新指令序列
\end{verbatim}
\end{quote}

\subsection{十、自测题与答案}

\subsubsection{题 1：寄存器重命名能消除哪类依赖？}

答案：能消除 Anti dependence 和 Output dependence，不能消除 True dependence。

\subsubsection{题 2：冲突图中一条边表示什么？}

答案：两个变量活跃区间重叠，不能放入同一个物理寄存器。

\subsubsection{题 3：局部寄存器分配中何时 spill？}

答案：某程序点活跃变量数超过可用寄存器数时要 spill。

\end{document}
